\chapter{Information Theory}
\section{Entropia}

L'entropia è un concetto fondamentale della teoria dell'informazione e permette di \textbf{quantificare la quantità media di informazione associata a una variabile aleatoria}.

Per comprenderne il significato, partiamo dal concetto di informazione associata all'osservazione di un singolo evento.

\subsection{Quantità di informazione di un evento}

Consideriamo una variabile aleatoria discreta $x$, descritta da una distribuzione di probabilità $p(x)$.

Vogliamo determinare quanta informazione otteniamo quando osserviamo che la variabile assume un determinato valore.

L'idea fondamentale è che \textbf{la quantità di informazione ricevuta dipende da quanto è sorprendente l'evento osservato}.

In particolare:
\begin{itemize}
    \item Un evento molto probabile fornisce poca informazione, perché ci aspettavamo che accadesse.
    \item Un evento poco probabile fornisce molta informazione, perché la sua osservazione è sorprendente.
    \item Un evento certo, cioè con probabilità $1$, non fornisce alcuna nuova informazione.
\end{itemize}

\paragraph{Esempio intuitivo.}

Supponiamo di ricevere una delle seguenti notizie:
\begin{itemize}
    \item \textbf{Evento A:} domani sorgerà il Sole, con probabilità praticamente uguale a $1$.
    \item \textbf{Evento B:} lanciando un dado equilibrato uscirà il numero $6$, con probabilità $1/6$.
\end{itemize}

La seconda notizia fornisce più informazione della prima, poiché descrive un evento meno prevedibile.

Cerchiamo quindi una funzione $h(x)$ che misuri l'informazione associata all'osservazione di $x$.

Questa funzione deve soddisfare due proprietà fondamentali:
\begin{enumerate}
    \item \textbf{Dipendenza dalla probabilità:} al diminuire di $p(x)$, la quantità di informazione $h(x)$ deve aumentare.
    
    \item \textbf{Additività per eventi indipendenti:} se due eventi $x$ e $y$ sono statisticamente indipendenti, l'informazione ottenuta osservandoli entrambi deve essere uguale alla somma delle informazioni ottenute separatamente:
    \[
        h(x,y)=h(x)+h(y).
    \]
\end{enumerate}

Ricordiamo che, per eventi indipendenti, vale:
\[
    p(x,y)=p(x)p(y).
\]

Pertanto, cerchiamo una funzione che trasformi il prodotto delle probabilità in una somma.

La funzione logaritmica soddisfa esattamente questa proprietà:
\[
    \log_2\bigl(p(x)p(y)\bigr)
    =
    \log_2 p(x)+\log_2 p(y).
\]

Introducendo un segno negativo, otteniamo la misura di informazione:

\[
    \boxed{h(x)=-\log_2 p(x)}
\]

Il segno negativo è necessario perché, essendo $0<p(x)\leq 1$, il logaritmo è minore o uguale a zero.

In questo modo:
\[
    h(x)\geq 0.
\]

La quantità $h(x)$ si misura in \textbf{bit}, poiché abbiamo utilizzato un logaritmo in base $2$.

\paragraph{Esempio numerico.}

Consideriamo tre eventi con probabilità differenti:
\[
\begin{aligned}
    p(x_1)&=1,\\
    p(x_2)&=\frac{1}{2},\\
    p(x_3)&=\frac{1}{8}.
\end{aligned}
\]

Calcoliamo la quantità di informazione associata a ciascun evento:

\[
\begin{aligned}
    h(x_1)&=-\log_2 1=0 \text{ bit},\\[4pt]
    h(x_2)&=-\log_2 \frac12=1 \text{ bit},\\[4pt]
    h(x_3)&=-\log_2 \frac18=3 \text{ bit}.
\end{aligned}
\]

Osserviamo che l'evento meno probabile produce la maggiore quantità di informazione.

\medskip
\noindent
\textbf{Attenzione:} $h(x)$ misura l'informazione ottenuta osservando uno specifico valore della variabile aleatoria. Non rappresenta ancora l'informazione media della distribuzione.

\subsection{Definizione di entropia}
Prima di parlare della definizione, tieni conto di una cosa.
Quando diciamo di "osservare" una variabile aleatoria, stiamo dicendo che, noi la variabile aleatoria l'abbiamo già definita, ad esempio in questo modo:
\[
\begin{aligned}
\Omega = \{T,C\} \\
X : \Omega \to S \subset \mathbb{R} \\
X \in \{0,1\} , X(T) = 0, X(C) = 1 \\
p_X(0) = P(X = 0) = 0.9 ; h(0) = 0.152 bit \\
p_X(1) = P(X = 1) = 0.1 h(1) = 3.322 bit  \\
\end{aligned}
\]
Una volta definita, possiamo trattare la variabile aleatoria, come un "filtro", ad esempio una volta definita, io faccio un esperimento (seguendo lo spazio di campionamento) e supponiamo esce "Testa", uscendo Testa
posso applicare la variabile aleatoria che mi restituisce 0, e la pmf, dice che la probabilità di uscire 0, era 0.9, l'informazione associata \textbf{a questo specifico evento è bassa} perchè è molto probabile che accada.

\begin{center}
    Quindi finora abbiamo visto la misura dell'informazione per un singolo evento, io adesso vorrei misurare l'informazione associata all'osservazione della variabile aleatoria.
    E quindi nelle successive pagine, andremo a vedere l'informazione media, che ci aspettiamo di ottenere osservando una variabile aleatoria.
\end{center}
Per farlo, consideriamo l'operatore di valore atteso per una variabile aleatoria:
\[
\mathbb{E}[X] = \sum_{x_i \in S} x_i P(X = x_i)
\]
è la somma di tutti i valori che la variabile può assumere, ciascuno pesato dalla probabilità con cui si verifica.
Ci dice in media, i valori che assume una variabile aleatoria.

Possiamo estendere questo concetto al valore atteso della funzione $h(x)$ rispetto alla distribuzione $p(x)$:

\[
    H[x]=\mathbb{E}_{x}[h(x)].
\]
Intuitivamente, l'entropia misura la quantità media di informazione che otteniamo osservando il valore assunto dalla variabile aleatoria \(x\), distribuita secondo la funzione di massa di probabilità \(p(x)\).

\newpage
Ricordando che il valore atteso di una funzione $f(x)$ è:

\[
    \mathbb{E}_{x}[f(x)]
    =
    \sum_x p(x)f(x),
\]

e sostituendo $h(x)=-\log_2 p(x)$, otteniamo:

\[
    H[x]
    =
    \sum_x p(x)\bigl(-\log_2 p(x)\bigr).
\]

Quindi:

\[
    \boxed{
        H[x]=-\sum_x p(x)\log_2 p(x)
    }
\]

Questa quantità prende il nome di \textbf{entropia della variabile aleatoria $x$}.

L'entropia misura quindi l'informazione media associata ai possibili valori della variabile, tenendo conto delle rispettive probabilità.

\paragraph{Interpretazione.}

È importante distinguere:
\begin{itemize}
    \item $h(x)$: informazione associata a una singola osservazione.
    \item $H[x]$: informazione media che ci aspettiamo di ottenere osservando la variabile aleatoria.
\end{itemize}
Ad esempio se
\[
H[x]=0.469\text{ bit }
\]
può essere considerata come l'informazione media, osservando quella variabile aleatoria, è una media.
Non significa che ogni osservazione fornisca esattamente \(0.469\) bit, ma che, considerando molte osservazioni della variabile aleatoria, la quantità di informazione per osservazione tende mediamente a quel valore.

L'entropia può essere interpretata anche come una misura dell'incertezza associata \textbf{all'osservazione della} variabile:
\begin{itemize}
    \item Se il risultato dell'osservazione è prevedibile, l'entropia è bassa.
    \item Se il risultato dell'osservazione è molto incerto, l'entropia è alta.
\end{itemize}

\begin{table}[h]
    \centering
    \begin{tabular}{|c|c|l|}
        \hline
        \textbf{Distribuzione} & \textbf{Entropia} & \textbf{Interpretazione} \\
        \hline
        $100\%$ Testa, $0\%$ Croce & $0$ bit & Risultato certo \\
        \hline
        $90\%$ Testa, $10\%$ Croce & $0.469$ bit & Risultato abbastanza prevedibile \\
        \hline
        $50\%$ Testa, $50\%$ Croce & $1$ bit & Massima incertezza \\
        \hline
    \end{tabular}
    \caption{Confronto dell'entropia per diverse distribuzioni di probabilità.}
    \label{tab:entropia_moneta}
\end{table}

\paragraph{Caso di probabilità nulla.}

Nella formula dell'entropia può comparire un valore $x$ tale che $p(x)=0$.

In questo caso, adottiamo la convenzione:

\[
    \boxed{0\log_2 0=0}
\]

La convenzione deriva dal limite:

\[
    \lim_{\varepsilon\to 0^+}
    \varepsilon\ln\varepsilon=0.
\]

Pertanto, gli eventi con probabilità nulla non contribuiscono al valore dell'entropia.
